\section{\texorpdfstring{Yoneda extensions,\allowbreak{} composition signs,\allowbreak{} splitting,\allowbreak{} and K-groups}{Yoneda  composition   and K-groups}}\label{proofs-8774-9196-yoneda-extensions-composition-signs-splitting-and-k-groups}

This is a private source-bound review of original Derived Categories lines 8774–\allowbreak{}9196,\allowbreak{} with a receiving correction in Cohomology. It continues PART12. The authority is Stacks commit a04446e5\allowbreak{}7ec1fbc2\allowbreak{}52a871af\allowbreak{}cec7752f\allowbreak{}b2807b14; the receiving public edition is 95a51593\allowbreak{}aceb247a\allowbreak{}c7a18311\allowbreak{}1678c3a0\allowbreak{}f8b8f4b5. The accompanying reading and operation records bind the exact bytes. This note proves the arguments used for review; it does not replace the translated source by its editorial strengthenings. No novelty is claimed.

Human source:\allowbreak{} \href{https://github.com/stacks/stacks-project/blob/a04446e57ec1fbc252a871afcec7752fb2807b14/derived.tex\#L8774}{The Stacks Project,\allowbreak{} Derived Categories,\allowbreak{} original source}. The \href{https://stacks.math.columbia.edu/tag/06XP}{current online Ext section} was checked as a secondary comparison on 2026-09-30; it does not replace the frozen authority. The arguments below use the original TeX,\allowbreak{} not a PDF.

\subsection{1. The exact pushout construction and common refinement}\label{proofs-8774-9196-the-exact-pushout-construction-and-common-refinement}

For a degree n extension E of B by A,\allowbreak{} put A in degree -n,\allowbreak{} its successive middle objects in degrees -n\allowbreak{}+1 through 0,\allowbreak{} and use its actual arrows as the differentials. Write C\_\allowbreak{}E for this complex,\allowbreak{} s\_\allowbreak{}E:\allowbreak{}C\_\allowbreak{}E\allowbreak{}→B[0]\allowbreak{} for the augmentation,\allowbreak{} and f\_\allowbreak{}E:\allowbreak{}C\_\allowbreak{}E\allowbreak{}→A[n]\allowbreak{} for the projection with component id\_\allowbreak{}A in degree -n. Thus δ(E)\allowbreak{}\allowbreak{}=f\_\allowbreak{}E s\_\allowbreak{}E\textasciicircum{}\{-1\}. The identity component is the natural projection in the source\textquotesingle s display; retaining it matters for Section 3.

Suppose s:\allowbreak{}L\allowbreak{}→B[0]\allowbreak{} is a quasi-isomorphism and f:\allowbreak{}L\allowbreak{}→A[n]\allowbreak{} is a chain map. Replace L by τ\_\allowbreak{}\{\textbackslash{}leq0\}L,\allowbreak{} restricting both maps. This keeps the represented morphism and makes L\textasciicircum{}r\allowbreak{}=0 for r\textgreater0. In particular H\textasciicircum{}r(L)\allowbreak{}\allowbreak{}=0 for r<0,\allowbreak{} and L\textasciicircum{}0\allowbreak{}→B is an epimorphism with kernel im(d\_\allowbreak{}L\textasciicircum{}\{-1\})\allowbreak{}. The component f\^{}\{-n\} kills im(d\_\allowbreak{}L\textasciicircum{}\{-n-1\})\allowbreak{}. Therefore it factors through Q\allowbreak{}=L\textasciicircum{}\{-n\}/\allowbreak{}im(d\_\allowbreak{}L\textasciicircum{}\{-n-1\})\allowbreak{}. Exactness at degree -n identifies Q with im(d\_\allowbreak{}L\textasciicircum{}\{-n\})\allowbreak{} and gives a monomorphism j:\allowbreak{}Q\allowbreak{}→L\textasciicircum{}\{-n\allowbreak{}+1\}.

Form the pushout

P\allowbreak{}=coker((d\_\allowbreak{}L\textasciicircum{}\{-n\},\allowbreak{}-f\textasciicircum{}\{-n\})\allowbreak{}:\allowbreak{}L\textasciicircum{}\{-n\}\allowbreak{}→L\textasciicircum{}\{-n\allowbreak{}+1\}\allowbreak{}⊕A)\allowbreak{}.

Equivalently,\allowbreak{} push j out along Q\allowbreak{}→A. The resulting map A\allowbreak{}→P,\allowbreak{} a\allowbreak{}↦[0,\allowbreak{}a]\allowbreak{},\allowbreak{} is a monomorphism,\allowbreak{} and coker(A\allowbreak{}→P)\allowbreak{}\allowbreak{}=coker(j)\allowbreak{}. This follows directly from the pushout of the exact sequence 0\allowbreak{}→Q\allowbreak{}→L\textasciicircum{}\{-n\allowbreak{}+1\}\allowbreak{}→coker(j)\allowbreak{}\allowbreak{}→0 in an abelian category. For n\textgreater1 the next map is [x,\allowbreak{}a]\allowbreak{}\allowbreak{}↦d\_\allowbreak{}L\textasciicircum{}\{-n\allowbreak{}+1\}x; it is well defined and its kernel is the image of A. The remaining terms are L\textasciicircum{}\{-n\allowbreak{}+2\},\allowbreak{}…,\allowbreak{}L\textasciicircum{}0 followed by B. For n\allowbreak{}=1 the next map is [x,\allowbreak{}a]\allowbreak{}\allowbreak{}↦s\textasciicircum{}0(x)\allowbreak{},\allowbreak{} and the same pushout sequence identifies its cokernel as zero and its kernel as A. Thus both cases give an exact degree n extension E(f,\allowbreak{}s)\allowbreak{}.

There is a chain map u:\allowbreak{}L\allowbreak{}→C\_\allowbreak{}\{E(f,\allowbreak{}s)\allowbreak{}\}. Its components are zero below -n,\allowbreak{} f\^{}\{-n\} in degree -n,\allowbreak{} x\allowbreak{}↦[x,\allowbreak{}0]\allowbreak{} in degree -n\allowbreak{}+1,\allowbreak{} and identities above that. The essential equation is {[}d x,\allowbreak{}0]\allowbreak{}\allowbreak{}=[0,\allowbreak{}f x]\allowbreak{},\allowbreak{} which is precisely the relation in the pushout. All other chain equations follow from f d\allowbreak{}=0 and d²\allowbreak{}=0. The equalities s\_\allowbreak{}E u\allowbreak{}=s and f\_\allowbreak{}E u\allowbreak{}=f are strict. On H\^{}0 the first equality shows that u is an isomorphism,\allowbreak{} and all other cohomology objects on either side vanish. Hence u is a quasi-isomorphism and δ(E(f,\allowbreak{}s)\allowbreak{})\allowbreak{}\allowbreak{}=f s\^{}\{-1\}. This proves the surjectivity in the source\textquotesingle s Yoneda classification.

For injectivity,\allowbreak{} let E,\allowbreak{}E' have equal classes. The common-denominator description of the localization of K(A)\allowbreak{} gives a complex L and quasi-isomorphisms t:\allowbreak{}L\allowbreak{}→C\_\allowbreak{}E and t':\allowbreak{}L\allowbreak{}→C\_\allowbreak{}\{E'\} such that s\_\allowbreak{}E t\allowbreak{}=s\_\allowbreak{}\{E'\}t' and f\_\allowbreak{}E t\allowbreak{}=f\_\allowbreak{}\{E'\}t' in K(A)\allowbreak{}. Replace L by τ\_\allowbreak{}\{\textbackslash{}leq0\}L. A homotopy from L to B[0]\allowbreak{} could contribute to degree 0 only through a map L\textasciicircum{}1\allowbreak{}→B. That term is zero. Consequently the first equality is already a strict equality of chain maps.

Set α\allowbreak{}=t\textasciicircum{}\{-n\}:\allowbreak{}L\textasciicircum{}\{-n\}\allowbreak{}→A and α'\allowbreak{}=t'\textasciicircum{}\{-n\}. The second equality says α-α'\allowbreak{}=h d\_\allowbreak{}L\textasciicircum{}\{-n\} for a map h:\allowbreak{}L\textasciicircum{}\{-n\allowbreak{}+1\}\allowbreak{}→A. There are no other possibly nonzero components in a homotopy to A[n]\allowbreak{}. Let ι':\allowbreak{}A\allowbreak{}→Z'\_\allowbreak{}\{n-1\} denote the first arrow of E\textquotesingle. Modify t\textquotesingle{} by the homotopy whose only component is h:\allowbreak{}

t̃'\textasciicircum{}\{-n\}\allowbreak{}=α,\allowbreak{} t̃'\textasciicircum{}\{-n\allowbreak{}+1\}\allowbreak{}=t'\textasciicircum{}\{-n\allowbreak{}+1\}\allowbreak{}+ι'h,\allowbreak{}

and keep every other component equal to t\textquotesingle. These formulas are t̃'\allowbreak{}=t'\allowbreak{}+d h\allowbreak{}+h d. They remain valid for n\allowbreak{}=1:\allowbreak{} the augmentation kills ι'h,\allowbreak{} so it is unchanged. Thus t̃\textquotesingle{} is a chain map,\allowbreak{} is homotopic to t',\allowbreak{} and has the same augmentation as t. Both maps now have the same top component α strictly.

Apply the pushout construction to (α,\allowbreak{}s\_\allowbreak{}E t)\allowbreak{}. The map from its first middle term to Z\_\allowbreak{}\{n-1\} is [x,\allowbreak{}a]\allowbreak{}\allowbreak{}↦t\textasciicircum{}\{-n\allowbreak{}+1\}x\allowbreak{}+ι a. It kills the relation (d x,\allowbreak{}-α x)\allowbreak{} by the chain equation for t. The analogous formula using t̃\textquotesingle{} maps to Z'\_\allowbreak{}\{n-1\}. The later components are those of t and t̃',\allowbreak{} and both endpoint maps are identities. This constructs the actual common middle extension in the definition,\allowbreak{} including the homotopy adjustment that the source leaves to the reader. Conversely,\allowbreak{} a map of extensions fixing A,\allowbreak{}B commutes with s and f and is a quasi-isomorphism,\allowbreak{} so it preserves δ. The classification follows,\allowbreak{} and the stated common-refinement relation is reflexive,\allowbreak{} symmetric and transitive because equality of δ is.

The omitted adjustment is supplied here as an editorial proof. The source explicitly omits details,\allowbreak{} so omission alone is not counted as another mathematical error.

\subsection{2. Naturality and the actual group law in degree one}\label{proofs-8774-9196-naturality-and-the-actual-group-law-in-degree-one}

For a map of extensions from E of B by A to F of B\textquotesingle{} by A\textquotesingle{} with endpoint maps β:\allowbreak{}B\allowbreak{}→B' and α:\allowbreak{}A\allowbreak{}→A',\allowbreak{} the induced chain map v satisfies s\_\allowbreak{}F v\allowbreak{}=βs\_\allowbreak{}E and f\_\allowbreak{}F v\allowbreak{}=α[n]\allowbreak{}f\_\allowbreak{}E. In D(A)\allowbreak{} this proves

δ(F)\allowbreak{}β\allowbreak{}=α[n]\allowbreak{}δ(E)\allowbreak{}.

Pushout of the first monomorphism along α and pullback of the last epimorphism along β exist in an abelian category and preserve the exact extension. The displayed identity proves their usual covariance and contravariance,\allowbreak{} with no assumption about enough injectives or projectives. Direct sums give δ(E\allowbreak{}⊕F)\allowbreak{}\allowbreak{}=δ(E)\allowbreak{}\allowbreak{}⊕δ(F)\allowbreak{}.

In degree one,\allowbreak{} morphisms of extensions fixing the endpoints are isomorphisms:\allowbreak{} apply the kernel and cokernel exact sequences to their commutative diagram with exact rows,\allowbreak{} or the short five lemma. Thus the common-refinement classification is the isomorphism classification in Homology\textquotesingle s definition. Its Baer sum is the direct sum followed by the pushout along Σ:\allowbreak{}A\allowbreak{}⊕A\allowbreak{}→A and pullback along Δ:\allowbreak{}B\allowbreak{}→B\allowbreak{}⊕B. Naturality gives

δ(E\allowbreak{}+F)\allowbreak{}\allowbreak{}=Σ\href{δ(E)⊕δ(F)}{1}Δ\allowbreak{}=δ(E)\allowbreak{}\allowbreak{}+δ(F)\allowbreak{}.

For the split extension A\allowbreak{}→A\allowbreak{}⊕B\allowbreak{}→B,\allowbreak{} projection A\allowbreak{}⊕B\allowbreak{}→A is a homotopy from f\_\allowbreak{}E to zero on its two-term complex,\allowbreak{} so δ(E)\allowbreak{}\allowbreak{}=0. Pushout by -id\_\allowbreak{}A gives additive inverses. This proves the group identification,\allowbreak{} not merely the set bijection. These are details supporting the original lemma.

One sign from an earlier source section must remain explicit. For a short exact sequence 0\allowbreak{}→A\allowbreak{}→Z\allowbreak{}→B\allowbreak{}→0,\allowbreak{} C\_\allowbreak{}E is the source cone of A[0]\allowbreak{}\allowbreak{}→Z[0]\allowbreak{}. Its cone differential is the positive inclusion A\allowbreak{}→Z. The source\textquotesingle s distinguished cone triangle has last map -p,\allowbreak{} so its connecting morphism is

∂\_\allowbreak{}E\allowbreak{}=-p s\_\allowbreak{}E\textasciicircum{}\{-1\}\allowbreak{}=-δ(E)\allowbreak{}.

This is the convention in current derived.tex 4286–\allowbreak{}4326. Both maps have been retained; the proof below never silently identifies them.

\subsection{3. Endpoint variance and the full signed splicing formula}\label{proofs-8774-9196-endpoint-variance-and-the-full-signed-splicing-formula}

Let E be a degree i extension of B by A,\allowbreak{} and E\textquotesingle{} a degree j extension of C by B. Their classes belong to Ext\textasciicircum{}i(B,\allowbreak{}A)\allowbreak{} and Ext\textasciicircum{}j(C,\allowbreak{}B)\allowbreak{}. The appropriate composition is δ(E)\allowbreak{}[j]\allowbreak{}δ(E')\allowbreak{}:\allowbreak{}C[0]\allowbreak{}\allowbreak{}→A[i\allowbreak{}+j]\allowbreak{}. Therefore the source display at 8913–\allowbreak{}8915 must have domain Ext\textasciicircum{}i(B,\allowbreak{}A)\allowbreak{}×Ext\textasciicircum{}j(C,\allowbreak{}B)\allowbreak{} and codomain Ext\textasciicircum{}\{i\allowbreak{}+j\}(C,\allowbreak{}A)\allowbreak{}. This is DERIVED-NEW-028. The preceding display with the general objects X,\allowbreak{}Y,\allowbreak{}Z is correct and is retained.

Let S be the splice with all the original arrows:\allowbreak{}

0\allowbreak{}→A\allowbreak{}→Z\_\allowbreak{}\{i-1\}\allowbreak{}→…\allowbreak{}→Z\_\allowbreak{}0\allowbreak{}→Z'\_\allowbreak{}\{j-1\}\allowbreak{}→…\allowbreak{}→Z'\_\allowbreak{}0\allowbreak{}→C\allowbreak{}→0.

The joining arrow is Z\_\allowbreak{}0\allowbreak{}→B\allowbreak{}→Z'\_\allowbreak{}\{j-1\}. Its kernel and image are exactly the required adjacent image and kernel,\allowbreak{} so S is exact. Let J\allowbreak{}=C\_\allowbreak{}S,\allowbreak{} K\allowbreak{}=C\_\allowbreak{}E,\allowbreak{} and L\allowbreak{}=C\_\allowbreak{}\{E'\}. Define a:\allowbreak{}J\allowbreak{}→L as follows:\allowbreak{} it is the identity on the Z\textquotesingle{} terms in degrees -j\allowbreak{}+1,\allowbreak{}…,\allowbreak{}0,\allowbreak{} the augmentation Z\_\allowbreak{}0\allowbreak{}→B in degree -j,\allowbreak{} and zero in all lower degrees. It is a chain map,\allowbreak{} its augmentation is s\_\allowbreak{}S,\allowbreak{} and it induces the identity on the only nonzero cohomology,\allowbreak{} C in degree 0. Hence a is a quasi-isomorphism.

The shift convention is d\_\allowbreak{}\{K[j]\allowbreak{}\}\textasciicircum{}r\allowbreak{}=(-1)\allowbreak{}\textasciicircum{}j d\_\allowbreak{}K\textasciicircum{}\{r\allowbreak{}+j\}. Define b:\allowbreak{}J\allowbreak{}→K[j]\allowbreak{} by

b\textasciicircum{}r\allowbreak{}=(-1)\allowbreak{}\textasciicircum{}\{j(r\allowbreak{}+j)\allowbreak{}\} id\_\allowbreak{}\{K\textasciicircum{}\{r\allowbreak{}+j\}\} (-i-j≤r≤-j)\allowbreak{},\allowbreak{} b\textasciicircum{}r\allowbreak{}=0 (r>-j)\allowbreak{}.

Outside the displayed support its components are zero. For r\textless-j the chain equation is the scalar identity

(-1)\allowbreak{}\textasciicircum{}j(-1)\allowbreak{}\textasciicircum{}\{j(r\allowbreak{}+j)\allowbreak{}\}\allowbreak{}=(-1)\allowbreak{}\textasciicircum{}\{j(r\allowbreak{}+1\allowbreak{}+j)\allowbreak{}\}.

At r\allowbreak{}=-j both sides of the chain equation are zero. Thus b is a chain map. Its degree -j component is id\_\allowbreak{}\{Z\_\allowbreak{}0\},\allowbreak{} and consequently

s\_\allowbreak{}E[j]\allowbreak{} b\allowbreak{}=f\_\allowbreak{}\{E'\}a.

In the derived category this equality and the invertibility of s\_\allowbreak{}E[j]\allowbreak{} compute the desired composition without omitting any shift:\allowbreak{}

δ(E)\allowbreak{}[j]\allowbreak{}δ(E')\allowbreak{} \allowbreak{}= f\_\allowbreak{}E[j]\allowbreak{} b(s\_\allowbreak{}\{E'\}a)\allowbreak{}\textasciicircum{}\{-1\} \allowbreak{}= (-1)\allowbreak{}\textasciicircum{}\{ij\} f\_\allowbreak{}S s\_\allowbreak{}S\textasciicircum{}\{-1\} \allowbreak{}= (-1)\allowbreak{}\textasciicircum{}\{ij\}δ(S)\allowbreak{}. (3.1)\allowbreak{}

The sign comes from b\textasciicircum{}\{-i-j\}\allowbreak{}=(-1)\allowbreak{}\textasciicircum{}\{ij\}id\_\allowbreak{}A. In particular,\allowbreak{} the natural projection convention in Section 1 does not send unsigned splicing to composition. DERIVED-NEW-029 records this sign and its propagation into the higher-vanishing proof. The correction makes f\_\allowbreak{}E\textasciicircum{}\{-i\}\allowbreak{}=id\_\allowbreak{}A explicit; it does not change the differentials or rescale the existing class. Because the source leaves the component of f implicit,\allowbreak{} the finding is also a convention-clarity issue:\allowbreak{} a different,\allowbreak{} degree-dependent choice of f would be a different convention and cannot silently repair this calculation.

Here is a counterexample to dropping the sign. Put R\allowbreak{}=Q[ε]\allowbreak{}/\allowbreak{}(ε²)\allowbreak{} and k\allowbreak{}=R/\allowbreak{}(ε)\allowbreak{}. The sequence E:\allowbreak{}0\allowbreak{}→k\allowbreak{}→R\allowbreak{}→k\allowbreak{}→0 has injection 1\allowbreak{}↦ε and quotient q. The complex P with P\textasciicircum{}r\allowbreak{}=R for r≤0,\allowbreak{} P\textasciicircum{}r\allowbreak{}=0 for r>0,\allowbreak{} differential ε in every negative degree,\allowbreak{} and augmentation q is a projective resolution of k:\allowbreak{} ker(ε)\allowbreak{}\allowbreak{}=εR\allowbreak{}=im(ε)\allowbreak{},\allowbreak{} and coker(ε)\allowbreak{}\allowbreak{}=k. The chain map P\allowbreak{}→C\_\allowbreak{}E has components id\_\allowbreak{}R in degree 0,\allowbreak{} q in degree -1,\allowbreak{} and zero below. It represents δ(E)\allowbreak{} by q:\allowbreak{}P\textasciicircum{}\{-1\}\allowbreak{}→k.

A lift g:\allowbreak{}P\allowbreak{}→P[1]\allowbreak{} of that class has g\textasciicircum{}r\allowbreak{}=(-1)\allowbreak{}\textasciicircum{}\{r\allowbreak{}+1\}id\_\allowbreak{}R for r≤-1 and zero in degree 0. The chain equation -ε g\textasciicircum{}r\allowbreak{}=g\textasciicircum{}\{r\allowbreak{}+1\}ε holds. Then δ(E)\allowbreak{}[1]\allowbreak{}δ(E)\allowbreak{} is represented in degree -2 by q g\textasciicircum{}\{-2\}\allowbreak{}=-q. The splice S has middle arrow ε:\allowbreak{}R\allowbreak{}→R; a map P\allowbreak{}→C\_\allowbreak{}S has components id\_\allowbreak{}R in degrees 0 and -1 and q in degree -2,\allowbreak{} so δ(S)\allowbreak{} is represented by q. Every differential in Hom\_\allowbreak{}R(P,\allowbreak{}k)\allowbreak{} is zero because ε acts by zero on k. Thus Ext\textasciicircum{}2\_\allowbreak{}R(k,\allowbreak{}k)\allowbreak{}\allowbreak{}=k,\allowbreak{} and q and -q are distinct. This verifies an actual failure of unsigned composition,\allowbreak{} not just a formal discrepancy of chain maps.

The signs are associative:\allowbreak{} either bracketing of three splices contributes ij\allowbreak{}+ik\allowbreak{}+jk modulo 2. In particular no associativity axiom of the graded derived composition has been changed.

\subsection{\texorpdfstring{4. Zero products,\allowbreak{} the extension diagram,\allowbreak{} and the set of lifts}{4. Zero  the extension  and the set of lifts}}\label{proofs-8774-9196-zero-products-the-extension-diagram-and-the-set-of-lifts}

Write E:\allowbreak{}0\allowbreak{}→A\allowbreak{}→Z\allowbreak{}→B\allowbreak{}→0 and E':\allowbreak{}0\allowbreak{}→B\allowbreak{}→Z'\allowbreak{}→C\allowbreak{}→0,\allowbreak{} with inclusion u:\allowbreak{}B\allowbreak{}→Z' and quotient v:\allowbreak{}Z'\allowbreak{}→C. Applying Hom\_\allowbreak{}D(-,\allowbreak{}A[1]\allowbreak{})\allowbreak{} to the source\textquotesingle s triangle for E\textquotesingle{} gives the exact segment

Hom(B,\allowbreak{}A)\allowbreak{} --∂-\allowbreak{}-> Ext\textasciicircum{}1(C,\allowbreak{}A)\allowbreak{} --v*-\allowbreak{}-> Ext\textasciicircum{}1(Z',\allowbreak{}A)\allowbreak{} --u*-\allowbreak{}-> Ext\textasciicircum{}1(B,\allowbreak{}A)\allowbreak{} --boundary-\allowbreak{}-> Ext\textasciicircum{}2(C,\allowbreak{}A)\allowbreak{}.

For e\allowbreak{}=δ(E)\allowbreak{} the last boundary is e[1]\allowbreak{}∂\_\allowbreak{}\{E'\}\allowbreak{}=-e[1]\allowbreak{}δ(E')\allowbreak{}. Its vanishing is therefore equivalent to vanishing of the product in the original lemma,\allowbreak{} even though the sign must be kept in the exact formula. By exactness this holds precisely when there is w\allowbreak{}∈Ext\textasciicircum{}1(Z',\allowbreak{}A)\allowbreak{} with u*w\allowbreak{}=e. Section 2 represents w by 0\allowbreak{}→A\allowbreak{}→W\allowbreak{}→Z'\allowbreak{}→0. Its pullback to B is isomorphic,\allowbreak{} fixing A and B,\allowbreak{} to E. Choose that isomorphism and compose it with W×\_\allowbreak{}\{Z'\}B\allowbreak{}→W to obtain Z\allowbreak{}→W.

The induced square is a pullback,\allowbreak{} so Z\allowbreak{}→W is a monomorphism,\allowbreak{} and its quotient is C:\allowbreak{} compose W\allowbreak{}→Z'\allowbreak{}→C,\allowbreak{} whose kernel is the inverse image of B. The result is precisely the diagram in original 8941–\allowbreak{}8966,\allowbreak{} with all rows and columns exact. Conversely such a diagram identifies Z with W×\_\allowbreak{}\{Z'\}B by the short five lemma,\allowbreak{} so its middle row gives a lift of e and the product vanishes. This proves both directions of the omitted proof.

DERIVED-UNDER-017 records the extra information supplied by this argument. The set

F\_\allowbreak{}e\allowbreak{}=\{w\allowbreak{}∈Ext\textasciicircum{}1(Z',\allowbreak{}A)\allowbreak{}:\allowbreak{}u*w\allowbreak{}=e\}

is empty if the product is nonzero. If it is nonempty,\allowbreak{} addition makes it a torsor under

im(v*)\allowbreak{} ≅ Ext\textasciicircum{}1(C,\allowbreak{}A)\allowbreak{}/\allowbreak{}im(∂:\allowbreak{}Hom(B,\allowbreak{}A)\allowbreak{}\allowbreak{}→Ext\textasciicircum{}1(C,\allowbreak{}A)\allowbreak{})\allowbreak{}.

Indeed the difference of two lifts lies in ker(u*)\allowbreak{}\allowbreak{}=im(v*)\allowbreak{},\allowbreak{} every such element can be added to a lift,\allowbreak{} and the resulting action of im(v*)\allowbreak{} is free. Exactness identifies its kernel description as the stated quotient. These are isomorphism classes of middle-row extensions admitting a compatible top row; a choice of the isomorphism to that top row is additional data. The assertion is deliberately not a classification of diagrams with that choice fixed.

The earlier boundary is ∂(a)\allowbreak{}\allowbreak{}=-a[1]\allowbreak{}δ(E')\allowbreak{}; in particular its image equals the image of the pushout classes a[1]\allowbreak{}δ(E')\allowbreak{}. This keeps the sign while identifying the same quotient group. If that quotient is zero,\allowbreak{} an existing middle-row class is unique. This consequence is editorial,\allowbreak{} without a replacement of the source lemma.

\subsection{5. The receiving ideal-filtration diagram}\label{proofs-8774-9196-the-receiving-ideal-filtration-diagram}

Current cohomology.tex 14916–\allowbreak{}14934 cites the zero-product criterion. For every index i for which these ideal powers are defined,\allowbreak{} retain exactly

A\allowbreak{}=I\textasciicircum{}\{i\allowbreak{}+2\}/\allowbreak{}I\textasciicircum{}\{i\allowbreak{}+3\},\allowbreak{} B\allowbreak{}=I\textasciicircum{}\{i\allowbreak{}+1\}/\allowbreak{}I\textasciicircum{}\{i\allowbreak{}+2\},\allowbreak{} C\allowbreak{}=I\textasciicircum{}i/\allowbreak{}I\textasciicircum{}\{i\allowbreak{}+1\},\allowbreak{} Z\allowbreak{}=I\textasciicircum{}\{i\allowbreak{}+1\}/\allowbreak{}I\textasciicircum{}\{i\allowbreak{}+3\},\allowbreak{} Z'\allowbreak{}=I\textasciicircum{}i/\allowbreak{}I\textasciicircum{}\{i\allowbreak{}+2\},\allowbreak{} W\allowbreak{}=I\textasciicircum{}i/\allowbreak{}I\textasciicircum{}\{i\allowbreak{}+3\}.

All maps are induced by inclusions or quotient maps of the original ideal powers. The needed rows and columns are

0\allowbreak{}→A\allowbreak{}→Z\allowbreak{}→B\allowbreak{}→0,\allowbreak{} 0\allowbreak{}→A\allowbreak{}→W\allowbreak{}→Z'\allowbreak{}→0,\allowbreak{} 0\allowbreak{}→Z\allowbreak{}→W\allowbreak{}→C\allowbreak{}→0,\allowbreak{} 0\allowbreak{}→B\allowbreak{}→Z'\allowbreak{}→C\allowbreak{}→0.

They are exact because for subobjects J⊆K⊆L the sequence 0\allowbreak{}→K/\allowbreak{}J\allowbreak{}→L/\allowbreak{}J\allowbreak{}→L/\allowbreak{}K\allowbreak{}→0 is exact. Applied to each consecutive triple of the actual ideal powers,\allowbreak{} this also proves commutation of every square. Thus Section 4 proves the product is zero. Both source connecting maps b\^{}i and b\textasciicircum{}\{i\allowbreak{}+1\} are negatives of the corresponding degree-one Yoneda classes,\allowbreak{} so the two minus signs cancel in b\textasciicircum{}\{i\allowbreak{}+1\}[1]\allowbreak{}b\textasciicircum{}i. Applying M⊗\^{}L- to this zero composite still gives zero; applying cohomology therefore gives the zero composite of the Bockstein maps.

Original cohomology.tex 14425–\allowbreak{}14426 describes W as an extension in the opposite order from its quotient sequence. DERIVED-NEW-031 corrects that receiving sentence by displaying the actual middle-row sequence 0\allowbreak{}→I\textasciicircum{}\{i\allowbreak{}+2\}/\allowbreak{}I\textasciicircum{}\{i\allowbreak{}+3\}\allowbreak{}→I\textasciicircum{}i/\allowbreak{}I\textasciicircum{}\{i\allowbreak{}+3\}\allowbreak{}→I\textasciicircum{}i/\allowbreak{}I\textasciicircum{}\{i\allowbreak{}+2\}\allowbreak{}→0,\allowbreak{} whose pullback along B\allowbreak{}→Z' is the top row. This both corrects the reversed description and supplies exactly the hypothesis of the cited criterion. No ideal factor,\allowbreak{} exponent,\allowbreak{} or connecting map has been discarded.

The zero-class use in current adequate.tex 1216–\allowbreak{}1232 also survives (3.1)\allowbreak{}:\allowbreak{} multiplication of a class by a sign does not change whether it is zero. That bounded receiver check concerns the invocation of Yoneda classification,\allowbreak{} not acceptance of the remainder of that chapter\textquotesingle s proof.

\subsection{\texorpdfstring{6. Higher vanishing,\allowbreak{} including p\allowbreak{}=0 and p\allowbreak{}=1}{6. Higher  including  and }}\label{proofs-8774-9196-higher-vanishing-including-p0-and-p1}

For the original higher-vanishing lemma,\allowbreak{} if p\allowbreak{}=0 then Hom(A,\allowbreak{}A)\allowbreak{}\allowbreak{}=0 for every A,\allowbreak{} whence id\_\allowbreak{}A\allowbreak{}=0. Every A is a zero object and so is every complex up to its unique maps. All claimed Ext groups vanish. This case cannot use a positive-degree Yoneda extension of degree p.

Now let p≥1 and i\textgreater p. Represent ξ by the exact extension E as in Section 1 and set C\allowbreak{}=im(Z\_\allowbreak{}p\allowbreak{}→Z\_\allowbreak{}\{p-1\})\allowbreak{}. This is defined also when p\allowbreak{}=1. There are the two exact extensions printed by the source,\allowbreak{} E\textquotesingle{} of B by C of degree p and E\textquotesingle\textquotesingle{} of C by A of degree i-p. Equation (3.1)\allowbreak{} gives the complete equality

δ(E)\allowbreak{}\allowbreak{}=(-1)\allowbreak{}\textasciicircum{}\{p(i-p)\allowbreak{}\}δ(E'')\allowbreak{}[p]\allowbreak{}δ(E')\allowbreak{}.

The assumption makes δ(E')\allowbreak{} zero,\allowbreak{} so ξ\allowbreak{}=0. The case i\allowbreak{}=p is the assumption itself. DERIVED-NEW-030 inserts the p\allowbreak{}=0 argument and uses the image formula for C:\allowbreak{} the source\textquotesingle s additional kernel formula contains Z\_\allowbreak{}\{-1\} when p\allowbreak{}=1 and neither formula is defined as written when p\allowbreak{}=0. The signed equality is the propagation of NEW-029,\allowbreak{} not another counted finding. The theorem itself remains valid with its original hypotheses.

\subsection{7. A smaller splitting hypothesis and all choices of splitting}\label{proofs-8774-9196-a-smaller-splitting-hypothesis-and-all-choices-of-splitting}

Let K\allowbreak{}∈D\textasciicircum{}b(A)\allowbreak{},\allowbreak{} and write H\_\allowbreak{}r\allowbreak{}=H\textasciicircum{}r(K)\allowbreak{}. Assume only

Ext\textasciicircum{}\{r-s\allowbreak{}+1\}(H\_\allowbreak{}r,\allowbreak{}H\_\allowbreak{}s)\allowbreak{}\allowbreak{}=0 for every r\textgreater s. (7.1)\allowbreak{}

This is weaker than requiring every degree at least two for every such pair,\allowbreak{} as the original lemma does. Choose a≤b containing the cohomology support. The canonical truncation triangle at r is

τ\_\allowbreak{}\{\textbackslash{}leq r-1\}K\allowbreak{}→τ\_\allowbreak{}\{\textbackslash{}leq r\}K\allowbreak{}→H\_\allowbreak{}r[-r]\allowbreak{}\allowbreak{}→(τ\_\allowbreak{}\{\textbackslash{}leq r-1\}K)\allowbreak{}[1]\allowbreak{}.

Inductively the lower truncation has an isomorphism τ\_\allowbreak{}\{\textbackslash{}leq r-1\}K≅\allowbreak{}⊕\_\allowbreak{}\{a≤s<r\}H\_\allowbreak{}s[-s]\allowbreak{} inducing identity on its cohomology. Consequently the group containing the last arrow is

\allowbreak{}⊕\_\allowbreak{}\{a≤s<r\} Hom\_\allowbreak{}D(H\_\allowbreak{}r[-r]\allowbreak{},\allowbreak{}H\_\allowbreak{}s[-s\allowbreak{}+1]\allowbreak{})\allowbreak{} \allowbreak{}= \allowbreak{}⊕\_\allowbreak{}\{a≤s<r\} Ext\textasciicircum{}\{r-s\allowbreak{}+1\}(H\_\allowbreak{}r,\allowbreak{}H\_\allowbreak{}s)\allowbreak{}\allowbreak{}=0.

The triangle splits. Choose the splitting that is compatible with its first and second maps; combined with the previous identity on cohomology,\allowbreak{} it gives τ\_\allowbreak{}\{\textbackslash{}leq r\}K≅\allowbreak{}⊕\_\allowbreak{}\{a≤s≤r\}H\_\allowbreak{}s[-s]\allowbreak{},\allowbreak{} still inducing the identity. The one-degree and zero-cohomology cases follow from canonical truncation quasi-isomorphisms. This completes the finite induction and proves (7.1)\allowbreak{} is sufficient. Without the group vanishing,\allowbreak{} the actual last arrow is the exact obstruction:\allowbreak{} its being zero is necessary and sufficient for this stage to split,\allowbreak{} whether or not its ambient group is zero.

Set T\allowbreak{}=\allowbreak{}⊕\_\allowbreak{}\{a≤r≤b\}H\_\allowbreak{}r[-r]\allowbreak{}. A matrix entry from H\_\allowbreak{}r[-r]\allowbreak{} to H\_\allowbreak{}s[-s]\allowbreak{} is Ext\textasciicircum{}\{r-s\}(H\_\allowbreak{}r,\allowbreak{}H\_\allowbreak{}s)\allowbreak{}. It vanishes for r\textless s by negative Ext vanishing. For r\allowbreak{}=s it is Hom(H\_\allowbreak{}r,\allowbreak{}H\_\allowbreak{}r)\allowbreak{},\allowbreak{} and induces precisely that map on H\^{}r. Thus endomorphisms of T inducing zero on every cohomology object form the strictly triangular ideal

N\allowbreak{}=\allowbreak{}⊕\_\allowbreak{}\{a≤s<r≤b\}Ext\textasciicircum{}\{r-s\}(H\_\allowbreak{}r,\allowbreak{}H\_\allowbreak{}s)\allowbreak{} ⊆ End\_\allowbreak{}D(T)\allowbreak{}.

Its multiplication is actual shifted composition. A product of m entries requires a strictly descending chain of m cohomological degrees; hence N\textasciicircum{}\{b-a\allowbreak{}+1\}\allowbreak{}=0. Every 1\allowbreak{}+n is invertible,\allowbreak{} with inverse 1-n\allowbreak{}+n²-…\allowbreak{}+(-1)\allowbreak{}\textasciicircum{}\{b-a\}n\textasciicircum{}\{b-a\}. Conversely an automorphism inducing identity on cohomology has the form 1\allowbreak{}+n. Once an identity-on-cohomology splitting T\allowbreak{}→K exists,\allowbreak{} all such splittings form a torsor under 1\allowbreak{}+N by precomposition. The action is free and transitive by composing one splitting with the inverse of another.

Therefore the extra vanishing Ext\textasciicircum{}\{r-s\}(H\_\allowbreak{}r,\allowbreak{}H\_\allowbreak{}s)\allowbreak{}\allowbreak{}=0 for r\textgreater s makes the identity-on-cohomology splitting unique. Without it there is no general claim of a canonical splitting. DERIVED-UNDER-018 comprises (7.1)\allowbreak{},\allowbreak{} the precise obstruction,\allowbreak{} and this description of the choices. All are proved consequences of the source argument,\allowbreak{} kept in editorial material.

The source\textquotesingle s global Ext²\allowbreak{}=0 conclusion follows from Section 6 and then (7.1)\allowbreak{}. The two checked More on Algebra receivers,\allowbreak{} current 21522–\allowbreak{}21531 and 28714–\allowbreak{}28725,\allowbreak{} assume projective or free cohomology. Their positive Ext groups vanish,\allowbreak{} so (7.1)\allowbreak{} and the uniqueness condition both hold. Their stated existence conclusions remain valid; no receiving body replacement is needed for the stronger editorial result.

\subsection{8. K-groups and all omitted verifications in this interval}\label{proofs-8774-9196-k-groups-and-all-omitted-verifications-in-this-interval}

Use the source\textquotesingle s convention that object collections are sets (or work in the chosen universe)\allowbreak{}. In a triangulated category,\allowbreak{} the triangle of zeros implies [0]\allowbreak{}\allowbreak{}=0. An isomorphism X\allowbreak{}→Y has a triangle with third object zero,\allowbreak{} so [X]\allowbreak{}\allowbreak{}=[Y]\allowbreak{}. The triangle X\allowbreak{}→0\allowbreak{}→X[1]\allowbreak{} gives [X[1]\allowbreak{}]\allowbreak{}\allowbreak{}=-[X]\allowbreak{}; induction in both directions gives [X[n]\allowbreak{}]\allowbreak{}\allowbreak{}=(-1)\allowbreak{}\textasciicircum{}n[X]\allowbreak{} for every integer n. Thus the definition using all objects has the required invariance under isomorphism. The report P02-E0068/\allowbreak{}OCC-01538 corrects the preposition in original 9067. There is no prior admitted correction of this locus.

For a finite exact sequence 0\allowbreak{}→M\_\allowbreak{}a\allowbreak{}→M\_\allowbreak{}\{a\allowbreak{}+1\}\allowbreak{}→…\allowbreak{}→M\_\allowbreak{}b\allowbreak{}→0 in an abelian category,\allowbreak{} let B\_\allowbreak{}r be the image entering M\_\allowbreak{}r. The short exact sequences 0\allowbreak{}→B\_\allowbreak{}r\allowbreak{}→M\_\allowbreak{}r\allowbreak{}→B\_\allowbreak{}\{r\allowbreak{}+1\}\allowbreak{}→0 give [M\_\allowbreak{}r]\allowbreak{}\allowbreak{}=[B\_\allowbreak{}r]\allowbreak{}\allowbreak{}+[B\_\allowbreak{}\{r\allowbreak{}+1\}]\allowbreak{}. The alternating sum telescopes,\allowbreak{} since B\_\allowbreak{}a\allowbreak{}=B\_\allowbreak{}\{b\allowbreak{}+1\}\allowbreak{}=0. Consequently any bounded long exact cohomology sequence of a triangle X\allowbreak{}→Y\allowbreak{}→Z gives χ(Y)\allowbreak{}\allowbreak{}=χ(X)\allowbreak{}\allowbreak{}+χ(Z)\allowbreak{},\allowbreak{} with χ(X)\allowbreak{}\allowbreak{}=Σ\_\allowbreak{}r(-1)\allowbreak{}\textasciicircum{}r[H\textasciicircum{}r(X)\allowbreak{}]\allowbreak{}. The signs follow by arranging its terms H\textasciicircum{}r(X)\allowbreak{},\allowbreak{}H\textasciicircum{}r(Y)\allowbreak{},\allowbreak{}H\textasciicircum{}r(Z)\allowbreak{},\allowbreak{}H\textasciicircum{}\{r\allowbreak{}+1\}(X)\allowbreak{},\allowbreak{}… in their actual sequence order:\allowbreak{} up to an overall sign the three coefficients in degree r are (-1)\allowbreak{}\textasciicircum{}r,\allowbreak{}-(-1)\allowbreak{}\textasciicircum{}r,\allowbreak{}(-1)\allowbreak{}\textasciicircum{}r.

The assignment A\allowbreak{}↦A[0]\allowbreak{} sends each short exact sequence to the source\textquotesingle s distinguished triangle and so induces K\_\allowbreak{}0(A)\allowbreak{}\allowbreak{}→K\_\allowbreak{}0(D\textasciicircum{}b(A)\allowbreak{})\allowbreak{}. The Euler assignment above descends in the reverse direction and their composition on [A]\allowbreak{} is the identity. For any K\allowbreak{}∈D\textasciicircum{}b(A)\allowbreak{},\allowbreak{} the finite canonical truncation triangles give

[K]\allowbreak{}\allowbreak{}=Σ\_\allowbreak{}r[H\textasciicircum{}r(K)\allowbreak{}[-r]\allowbreak{}]\allowbreak{}\allowbreak{}=Σ\_\allowbreak{}r(-1)\allowbreak{}\textasciicircum{}r[H\textasciicircum{}r(K)\allowbreak{}[0]\allowbreak{}]\allowbreak{}.

The finite cohomological support makes this induction terminate at zero. This proves the other composition is the identity,\allowbreak{} without assuming the initial representative of K is termwise bounded. If needed,\allowbreak{} a bounded representative is obtained by the actual canonical truncations at the two cohomological endpoints. The source\textquotesingle s alternative finite induction on stupid truncations of that representative is valid.

If B⊆A is weak Serre,\allowbreak{} it is a full abelian subcategory whose inclusion is exact. In the long cohomology sequence of a triangle in D\textasciicircum{}b\_\allowbreak{}B(A)\allowbreak{},\allowbreak{} all terms lie in B. Kernels and cokernels of its maps belong to B,\allowbreak{} so all intervening images belong to B as well. The preceding telescoping argument therefore takes place in K\_\allowbreak{}0(B)\allowbreak{},\allowbreak{} not just K\_\allowbreak{}0(A)\allowbreak{}. The canonical truncations stay in D\textasciicircum{}b\_\allowbreak{}B(A)\allowbreak{},\allowbreak{} since their cohomology objects are retained or replaced by zero. The same formula gives mutually inverse maps K\_\allowbreak{}0(B)\allowbreak{}⇄K\_\allowbreak{}0(D\textasciicircum{}b\_\allowbreak{}B(A)\allowbreak{})\allowbreak{}. This fills the proof at original 9164–\allowbreak{}9176. It does not assume that every such complex has a representative with all its terms in B.

An exact functor of triangulated categories sends each triangle relation to a triangle relation,\allowbreak{} so its assignment on objects descends to K\_\allowbreak{}0. For the source\textquotesingle s homological functor H with finite support on every shift orbit,\allowbreak{} apply its long exact sequence to a triangle. The union of the three finite supports is finite,\allowbreak{} and the same Euler calculation proves that [X]\allowbreak{}\allowbreak{}↦Σ\_\allowbreak{}r(-1)\allowbreak{}\textasciicircum{}r[H(X[r]\allowbreak{})\allowbreak{}]\allowbreak{} respects every relation. Finally,\allowbreak{} a bifunctor exact in each variable respects the triangle relations separately in each input. Its map on pairs of generators therefore descends to the claimed bilinear map on the two K-groups. These arguments complete the omitted verifications through original 9196.

\subsection{\texorpdfstring{9. Dispositions,\allowbreak{} propagation,\allowbreak{} and limits}{9.   and limits}}\label{proofs-8774-9196-dispositions-propagation-and-limits}

One additional physical report is complete:\allowbreak{} P02-E0068/\allowbreak{}OCC-01538,\allowbreak{} DERIVED-RECON-069,\allowbreak{} missing copyedit. Cumulatively that is 117 of 176 reports,\allowbreak{} 69 groups,\allowbreak{} 108 already corrected,\allowbreak{} eight missing copyedit/\allowbreak{}notation reports and one missing mathematical-index report.

New unreported findings are NEW-028 (Ext endpoints)\allowbreak{},\allowbreak{} NEW-029 (signed splicing with its actual projection convention)\allowbreak{},\allowbreak{} NEW-030 (p\allowbreak{}=0/\allowbreak{}p\allowbreak{}=1 proof endpoints)\allowbreak{},\allowbreak{} and NEW-031 (receiving Cohomology extension sequence)\allowbreak{}. There are eight new proposed operations:\allowbreak{} seven in Derived Categories and one in Cohomology. The first seven include the received copyedit. The two new editorial consequences are UNDER-017 and UNDER-018.

Propagation has been proved into the higher-vanishing lemma,\allowbreak{} the zero-product criterion,\allowbreak{} the ideal-filtration Bockstein argument,\allowbreak{} the bounded splitting lemmas,\allowbreak{} and the bounded receiving uses specified above. The common-refinement and Baer-sum arguments support correct source claims and are not counted as errors merely because their source proofs are short. The K-group section requires only its received preposition correction.

Semantic review is complete through 9196; sequential reading reaches 9232. The next section is Unbounded complexes,\allowbreak{} beginning at 9206; the first incomplete source lemma begins at 9214. New unread source starts at 9233. No claim is made about that lemma\textquotesingle s proof yet. The Illusie I.3.0–\allowbreak{}I.3.1.1 producer relay received during this review is queued separately,\allowbreak{} without mathematical acceptance or interruption of the correction priority. No new sessions,\allowbreak{} live source composition,\allowbreak{} builds,\allowbreak{} Lean runs,\allowbreak{} publication,\allowbreak{} or cleanup were performed.
